\section*{Bab \ref{ch:g12:seq} --- Barisan}

\begin{solution}{exo:g12:seq:1}
Misalkan $P(n)$ pernyataan $\sum_{k=1}^{n} k^2 = \frac{n(n+1)(2n+1)}{6}$.
\emph{Langkah dasar:} untuk $n = 0$ kedua ruasnya $0$ (jumlah kosong).
\emph{Langkah induksi:} andaikan $P(n)$. Maka
\[
\sum_{k=1}^{n+1} k^2 = \frac{n(n+1)(2n+1)}{6} + (n+1)^2
= \frac{(n+1)\bigl(n(2n+1) + 6(n+1)\bigr)}{6}
= \frac{(n+1)(2n^2 + 7n + 6)}{6}.
\]
Karena $2n^2 + 7n + 6 = (n+2)(2n+3)$, bentuk ini sama dengan
$\frac{(n+1)(n+2)(2(n+1)+1)}{6}$, yaitu $P(n+1)$. Menurut induksi, $P(n)$
berlaku untuk semua $n$.
\end{solution}

\begin{solution}{exo:g12:seq:2}
$a_{n+1} - a_n = \frac{n+2}{n+1} - \frac{n+1}{n}
= \frac{n(n+2) - (n+1)^2}{n(n+1)} = \frac{-1}{n(n+1)} < 0$: jadi $(a_n)$
tegas turun.

$(b_n)$ mempunyai suku positif dan
$\frac{b_{n+1}}{b_n} = \frac{2^{n+1}}{n+1}\cdot\frac{n}{2^n} =
\frac{2n}{n+1} \geq 1 \iff 2n \geq n+1 \iff n \geq 1$: jadi $(b_n)$ \omterm{def:g12:seq:monotonic}{naik}
(tegas untuk $n \geq 2$).

$c_{n+1} - c_n = (n+1)^2 - 10(n+1) - n^2 + 10n = 2n - 9$, yang bernilai negatif
untuk $n \leq 4$ dan positif untuk $n \geq 5$: jadi $(c_n)$ turun sampai
$c_5 = -25$, yaitu \omterm{def:g10:functions:extrema}{minimumnya}, lalu \omterm{def:g12:seq:monotonic}{naik}. \omterm{def:g12:seq:sequence}{Barisan} itu tidak \omterm{def:g12:seq:monotonic}{monoton}.
\end{solution}

\begin{solution}{exo:g12:seq:3}
Keluarkan suku yang menguasainya sebagai faktor:
\[
u_n = \frac{n^2(3 - 1/n + 1/n^2)}{n^2(2 + 5/n^2)} \xrightarrow[n\to+\infty]{} \frac{3}{2}.
\]
Kalikan dengan bentuk sekawannya:
\[
v_n = \frac{(n+1) - n}{\sqrt{n+1} + \sqrt{n}} = \frac{1}{\sqrt{n+1}+\sqrt{n}}
\xrightarrow[n\to+\infty]{} 0.
\]
Bagilah pembilang dan penyebutnya dengan $3^n$:
\[
w_n = \frac{(2/3)^n - 1}{1 + (1/3)^n} \xrightarrow[n\to+\infty]{} \frac{0-1}{1+0} = -1,
\]
dengan memakai $\lim q^n = 0$ untuk $\abs{q} < 1$.
\end{solution}

\begin{solution}{exo:g12:seq:4}
$u_n = 5 + 3n \to +\infty$ dan $v_n = 8 \cdot (1/2)^n = 2^{3-n} \to 0$.
Jumlahnya adalah
\[
\sum_{k=0}^{n} u_k = (n+1)\,\frac{5 + (5+3n)}{2} = \frac{(n+1)(10+3n)}{2}
\xrightarrow[n\to+\infty]{} +\infty,
\]
\[
\sum_{k=0}^{n} v_k = 8\,\frac{1 - (1/2)^{n+1}}{1 - 1/2}
= 16\left(1 - \left(\tfrac{1}{2}\right)^{n+1}\right)
\xrightarrow[n\to+\infty]{} 16 .
\]
\end{solution}

\begin{solution}{exo:g12:seq:5}
Karena $-1 \leq \cos n \leq 1$,
\[
\frac{n-1}{n+1} \leq \frac{n + \cos n}{n+1} \leq 1,
\]
dan $\frac{n-1}{n+1} \to 1$, sehingga limitnya $1$ menurut teorema apit.

Untuk limit yang kedua, tulislah
\[
0 \leq \frac{n!}{n^n}
= \frac{1}{n}\cdot\frac{2}{n}\cdots\frac{n}{n}
\leq \frac{1}{n},
\]
sebab setiap faktor $\frac{k}{n}$ dengan $2 \leq k \leq n$ paling besar $1$.
Karena $\frac1n \to 0$, teorema apit memberi $\frac{n!}{n^n} \to 0$. (Batas
\omterm{def:g12:seq:arith-geom}{geometri} yang disarankan juga berhasil: setiap faktor dengan $k \leq n/2$
paling besar $\frac12$, yang memberi batas lebih kuat $(1/2)^{\floor{n/2}}$.)
\end{solution}

\begin{solution}{exo:g12:seq:6}
\emph{1.} $u_0 = 0 \in \intcc{0}{2}$. Jika $0 \leq u_n \leq 2$, maka
$2 \leq u_n + 2 \leq 4$, sehingga $\sqrt{2} \leq u_{n+1} \leq 2$; khususnya
$0 \leq u_{n+1} \leq 2$. Menurut induksi, sifat itu berlaku untuk semua $n$.

\emph{2.} $u_{n+1} - u_n = \sqrt{u_n + 2} - u_n$. Untuk $x \in \intcc{0}{2}$,
$\sqrt{x+2} \geq x \iff x + 2 \geq x^2 \iff (2-x)(x+1) \geq 0$, yang benar.
Jadi $(u_n)$ \omterm{def:g12:seq:monotonic}{naik}.

\emph{3.} Karena \omterm{def:g12:seq:monotonic}{naik} dan \omterm{def:g12:seq:bounded}{terbatas di atas} oleh $2$, $(u_n)$ \omterm{def:g12:seq:limit}{konvergen} ke suatu
$\ell \in \intcc{0}{2}$. Menerapkan limit pada $u_{n+1} = \sqrt{u_n + 2}$
(pemetaan $x \mapsto \sqrt{x+2}$ kontinu) memberi $\ell = \sqrt{\ell + 2}$,
sehingga $\ell^2 - \ell - 2 = 0$, yaitu $\ell \in \{-1, 2\}$. Karena
$\ell \geq 0$, diperoleh $\lim u_n = 2$.
\end{solution}

\begin{solution}{exo:g12:seq:7}
\emph{1.} Di antara dua dosis, $40\%$ obatnya tersingkirkan, sehingga
banyaknya $u_n$ menjadi $0.6\,u_n$; dosis berikutnya menambahkan $1$ satuan:
$u_{n+1} = 0.6\,u_n + 1$.

\emph{2.} $v_{n+1} = u_{n+1} - 2.5 = 0.6\,u_n + 1 - 2.5 = 0.6(u_n - 2.5)
= 0.6\,v_n$: jadi $(v_n)$ \omterm{def:g12:seq:arith-geom}{geometri} dengan \omterm{def:g11:seq:geometric}{rasio} $0.6$ dan suku pertama
$v_0 = 1 - 2.5 = -1.5$. Karena itu $v_n = -1.5 \times 0.6^n$ dan
\[
u_n = 2.5 - 1.5 \times 0.6^n .
\]

\emph{3.} Karena $0.6^n \to 0$, diperoleh $u_n \to 2.5$: jadi banyaknya obat
menetap pada $2.5$ satuan.
\end{solution}

\begin{solution}{exo:g12:seq:8}
\emph{1.} $u_0 = 3 > 1$. Jika $u_n > 1$, maka $u_n + 2 > 0$ dan
\[
u_{n+1} - 1 = \frac{4u_n - 1 - u_n - 2}{u_n + 2} = \frac{3(u_n - 1)}{u_n + 2} > 0 .
\]
Menurut induksi, $u_n > 1$ untuk semua $n$ (dan khususnya $u_n + 2 \neq 0$,
sehingga \omterm{def:g12:seq:sequence}{barisannya} terdefinisi dengan baik).

\emph{2.} Dengan memakai identitas di atas,
\[
v_{n+1} = \frac{1}{u_{n+1} - 1} = \frac{u_n + 2}{3(u_n - 1)}
= \frac{(u_n - 1) + 3}{3(u_n - 1)} = \frac{1}{3} + v_n .
\]
Jadi $(v_n)$ \omterm{def:g12:seq:arith-geom}{aritmetika} dengan \omterm{def:g11:seq:arithmetic}{beda} $\frac13$ dan
$v_0 = \frac{1}{u_0 - 1} = \frac12$.

\emph{3.} $v_n = \frac12 + \frac{n}{3}$, sehingga
$u_n = 1 + \frac{1}{v_n} = 1 + \frac{6}{3 + 2n}$. Karena $v_n \to +\infty$,
diperoleh $u_n \to 1$.
\end{solution}

\begin{solution}{exo:g12:seq:9}
\emph{1.} $H_{2n} - H_n = \sum_{k=n+1}^{2n} \frac{1}{k}$ adalah jumlah $n$
suku, masing-masing paling sedikit $\frac{1}{2n}$; sehingga
$H_{2n} - H_n \geq n \cdot \frac{1}{2n} = \frac12$.

\emph{2.} Dengan induksi pada $k$: $H_{2^0} = H_1 = 1 \geq 1$. Jika
$H_{2^k} \geq 1 + \frac{k}{2}$, maka menerapkan butir 1 dengan $n = 2^k$,
\[
H_{2^{k+1}} \geq H_{2^k} + \frac12 \geq 1 + \frac{k+1}{2}.
\]
\omterm{def:g12:seq:sequence}{Barisan} $(H_n)$ \omterm{def:g12:seq:monotonic}{naik} (setiap langkahnya menambahkan $\frac{1}{n+1} > 0$) dan
subbarisan $H_{2^k}$ takterbatas, sehingga $(H_n)$ tidak \omterm{def:g12:seq:bounded}{terbatas di atas}.
Karena \omterm{def:g12:seq:monotonic}{naik} dan takterbatas, \omterm{def:g12:seq:sequence}{barisan} itu menuju $+\infty$
(\cref{thm:g12:seq:monotone}).
\end{solution}

\begin{solution}{exo:g12:seq:10}
\emph{1.} \omterm{def:g12:seq:sequence}{Barisan} $d_n = b_n - a_n$ memenuhi
$d_{n+1} - d_n = (b_{n+1} - b_n) - (a_{n+1} - a_n) \leq 0$, sehingga $(d_n)$
turun; karena $d_n \to 0$, diperoleh $d_n \geq 0$ untuk semua $n$ (\omterm{def:g12:seq:sequence}{barisan}
turun yang mempunyai suku negatif akan tetap berada di bawahnya selamanya,
sehingga menghalangi limit $0$). Jadi $a_n \leq b_n$.

\emph{2.} Dari $a_n \leq b_n \leq b_0$, \omterm{def:g12:seq:sequence}{barisan} \omterm{def:g12:seq:monotonic}{naik} $(a_n)$
\omterm{def:g12:seq:bounded}{terbatas di atas}, sehingga \omterm{def:g12:seq:limit}{konvergen} ke suatu $\ell$. Serupa dengan itu
$(b_n)$, yang turun dan \omterm{def:g12:seq:bounded}{terbatas di bawah} oleh $a_0$, \omterm{def:g12:seq:limit}{konvergen} ke suatu
$\ell'$. Lalu $\ell' - \ell = \lim (b_n - a_n) = 0$, sehingga $\ell = \ell'$.

\emph{3.} $(a_n)$ bersifat (tegas) \omterm{def:g12:seq:monotonic}{naik} sebab
$a_{n+1} - a_n = \frac{1}{(n+1)!} > 0$. Untuk $(b_n)$,
\[
b_{n+1} - b_n = \frac{1}{(n+1)!} + \frac{1}{(n+1)(n+1)!} - \frac{1}{n\,n!}
= \frac{n(n+1) + n - (n+1)^2}{n(n+1)(n+1)!}
= \frac{-1}{n(n+1)(n+1)!} < 0 ,
\]
sehingga $(b_n)$ turun. Akhirnya $b_n - a_n = \frac{1}{n\,n!} \to 0$. Kedua
\omterm{def:g12:seq:sequence}{barisan} itu berdampingan, sehingga \omterm{def:g12:seq:limit}{konvergen} ke limit yang sama.
\end{solution}

\begin{solution}{pb:g12:seq:1}
\textbf{1.} Benar untuk $n = 1$ ($1 = 1^2$). Jika
$1 + 3 + \dots + (2n - 1) = n^2$, maka menambahkan bilangan ganjil
berikutnya: $n^2 + (2n + 1) = (n + 1)^2$: itulah keturunannya. Menurut
induksi, benar untuk semua $n \geq 1$.

\textbf{2.} $2^0 = 1 > 0$. Jika $2^n > n$, maka
$2^{n+1} = 2 \cdot 2^n > 2n \geq n + 1$ untuk $n \geq 1$ (dan
$n = 0$ terperiksa langsung): keturunannya beres.

\textbf{3.} $n = 0$: $1 \geq 1$. Jika
$(1 + x)^n \geq 1 + nx$, kalikan dengan $1 + x \geq 1 > 0$:
$(1 + x)^{n+1} \geq (1 + nx)(1 + x) = 1 + (n + 1)x + nx^2
\geq 1 + (n + 1)x$.

\textbf{4.} Langkah dari $n = 1$ ke $n = 2$: di antara dua
kelereng, ``$n$ yang pertama'' dan ``$n$ yang terakhir'' adalah dua
kelereng tunggal yang \emph{lepas} --- tidak ada kelereng bersama yang
menjembatani kedua kelompoknya, sehingga tidak ada yang memaksa warnanya
sama. Alasan keturunannya diam-diam menuntut kedua kelompoknya bertindih,
dan itu baru benar mulai $n \geq 2$; dengan langkah dasar $n = 1$ rantainya
tidak pernah bermula.

\textbf{5.} $4^0 - 1 = 0 = 3 \times 0$. Jika $4^n - 1 = 3k$,
maka $4^{n+1} - 1 = 4(4^n - 1) + 3 = 3(4k + 1)$: itulah keturunannya.

\textbf{6.} $x_1 = \frac32$, $x_2 = \frac{17}{12}$,
$x_3 = \frac{577}{408}$.

\textbf{7.} $x_{n+1}^2 - 2 = \frac{(x_n^2 + 2)^2 - 8x_n^2}
{4x_n^2} = \frac{(x_n^2 - 2)^2}{4 x_n^2}$: sebuah kuadrat dibagi
bilangan positif, jadi $\geq 0$, dan $> 0$ setiap kali $x_n^2 \neq 2$.
Induksinya: $x_0 = 2 > 0$ dengan $x_0^2 = 4 > 2$; jika $x_n > 0$
dan $x_n^2 > 2$, maka $x_{n+1}$ (rata-rata dua bilangan positif) bernilai
positif dan $x_{n+1}^2 - 2 > 0$.

\textbf{8.} $x_{n+1} - x_n = \frac{2 - x_n^2}{2x_n} < 0$ menurut
pertanyaan 7: jadi tegas turun.

\textbf{9.} Karena turun dan \omterm{def:g12:seq:bounded}{terbatas di bawah} (oleh $1$, sebab
$x_n^2 > 2 > 1$ dan $x_n > 0$): menurut teorema kekonvergenan
\omterm{def:g12:seq:monotonic}{monoton}, $(x_n)$ \omterm{def:g12:seq:limit}{konvergen} ke suatu $L \geq 1$.

\textbf{10.} Limit menghormati aljabarnya: dari
$x_{n+1} = \frac12\left(x_n + \frac{2}{x_n}\right)$ dan
$x_n \to L \geq 1 > 0$ diperoleh
$L = \frac12\left(L + \frac2L\right)$, sehingga $L^2 = 2$ dan, karena $L$
positif, $L = \sqrt2$. Vonisnya: kekonvergenannya terbukti,
limitnya terkenali --- Heron dibebaskan dengan pujian.

\textbf{11.} $x_{n+1} - \sqrt2 = \frac{x_n^2 - 2\sqrt2\,x_n +
2}{2x_n} = \frac{(x_n - \sqrt2)^2}{2x_n}$: tepat
$e_{n+1} = \frac{e_n^2}{2x_n}$, dan $x_n > \sqrt2$ memberi
$e_{n+1} \leq \frac{e_n^2}{2\sqrt2}$. Galat yang dikuadratkan: setiap
langkahnya melipatduakan banyaknya angka desimal yang benar, sebagaimana
teramati sejak jilid sebelumnya.

\textbf{12.} $e_0 \approx 0.5858$, $e_1 \approx 0.0858$,
$e_2 \approx 0.00245$, $e_3 \approx 2.1 \times 10^{-6}$.
\omterm{def:g11:seq:geometric}{Rasionya} $\frac{e_1}{e_0^2} \approx 0.25 = \frac{1}{2x_0}$;
$\frac{e_2}{e_1^2} \approx 0.333 = \frac{1}{2x_1}$;
$\frac{e_3}{e_2^2} \approx 0.353 \approx \frac{1}{2x_2}$: jadi
teoremanya sedang bekerja.

\textbf{13.} $H_{2^{18}} \geq 1 + 9 = 10$: sekitar
$260\,000$ suku ($2^{18} = 262\,144$) hanya untuk melampaui $10$ ---
kedivergenan yang merangkak (dan $H_n > 100$ akan menuntut suku lebih banyak
daripada atom di perpustakaan mana pun).

\textbf{14.} $S_n = 2 - \frac{1}{2^n}$ (jumlah \omterm{def:g12:seq:arith-geom}{geometri}), dan
$\frac{1}{2^n} \to 0$ (\cref{thm:g12:seq:geometric}), sehingga
$S_n \to 2$: batang cokelat yang digigit tanpa henti menuju,
tanpa pernah mencapai, keseluruhannya --- dan kini dalam bahasa
resmi limit.

\textbf{15.} Untuk $k \geq 2$:
$\frac{1}{k^2} \leq \frac{1}{k(k-1)} = \frac{1}{k-1} -
\frac1k$, sehingga
$S_n \leq 1 + \left(1 - \frac1n\right) < 2$: karena \omterm{def:g12:seq:monotonic}{naik} dan
\omterm{def:g12:seq:bounded}{terbatas di atas}, \omterm{def:g12:seq:sequence}{barisan} itu \omterm{def:g12:seq:limit}{konvergen} (kekonvergenan \omterm{def:g12:seq:monotonic}{monoton}). Euler
kemudian menamai limitnya: $\frac{\pi^2}{6}$.

\textbf{16.} Suku yang menuju $0$ memang \emph{perlu} agar
jumlahnya mengendap, tetapi tidak memutuskan apa-apa: suku harmonik
$\frac1n \to 0$ padahal jumlahnya meledak; suku
$\frac{1}{n^2} \to 0$ dan jumlahnya \omterm{def:g12:seq:limit}{konvergen}. Seberapa \emph{cepat}
sukunya mati, itulah seluruh pertanyaannya --- yaitu teori deret,
yang dibangun pada jilid universitas.

\textbf{17.} $a_1 = \sqrt2 \approx 1.41421$, $b_1 = 1.5$;
$a_2 \approx 1.45648$, $b_2 \approx 1.45711$: dua iterasi saja
sudah cocok sampai tiga angka desimal --- laju yang memukau.

\textbf{18.} $b_{n+1} - a_{n+1} = \frac{a_n + b_n}{2} -
\sqrt{a_n b_n} = \frac{(\sqrt{b_n} - \sqrt{a_n})^2}{2} \geq
0$: jadi kedua rataannya tetap berurutan. $(a_n)$ \omterm{def:g12:seq:monotonic}{naik}:
$a_{n+1} = \sqrt{a_n b_n} \geq \sqrt{a_n \cdot a_n} = a_n$;
$(b_n)$ turun secara setangkup.

\textbf{19.} $\frac{b_{n+1} - a_{n+1}}{b_n - a_n} =
\frac{(\sqrt{b_n} - \sqrt{a_n})^2}
{2(\sqrt{b_n} - \sqrt{a_n})(\sqrt{b_n} + \sqrt{a_n})}
= \frac{\sqrt{b_n} - \sqrt{a_n}}{2(\sqrt{b_n} + \sqrt{a_n})}
\leq \frac12$: jadi selisihnya paling sedikit terbagi dua, sehingga
$b_n - a_n \to 0$; bersama pertanyaan 18, kedua \omterm{def:g12:seq:sequence}{barisan} itu berdampingan dan
berbagi limit $M(1, 2)$.

\textbf{20.} Iterasi ketiganya memberi
$a_3 \approx b_3 \approx 1.456791$: jadi $M(1, 2) \approx
1.456791$ dalam tiga putaran engkol (selisihnya kira-kira
\emph{dikuadratkan}, seperti pada Heron). Urutan laju \omterm{def:g12:seq:sequence}{barisan} pada soal
ini, dari yang paling lambat sampai yang paling cepat: jumlah harmonik
(kedivergenan sedingin gletser), jumlah \omterm{def:g12:seq:arith-geom}{geometri} (galatnya terbagi dua setiap
langkah), lalu Heron dan rataan aritmetika--geometri (galatnya dikuadratkan
setiap langkah) --- dan laju yang tak wajar dari rataan itulah yang
memberitahu Gauss bahwa ia telah menemukan urat baru analisis.

\end{solution}
